\chapter{BV-вакуумный quotient полного 42-мерного носителя}
\label{ch:v8-full-42-carrier-bv-vacuum-quotient-gate}

Предыдущий гейт оставил открытыми два независимых вопроса: физическую
фермионную кратность и совместимость скалярного гессиана с
голдстоуновскими направлениями. Первый вопрос закрывается полной
хиральной проекцией. Для внутренней градуировки
\begin{equation}
 \Gamma_F=\operatorname{diag}(I_{11},-I_{10}),\qquad
 \{\Gamma_F,D_F\}=0,
 \quad
 P_{\rm phys}=\frac{I_{84}+\gamma_5\otimes\Gamma_F}{2},
 \quad \rank P_{\rm phys}=42.
 \label{eq:v8-full-carrier-physical-chiral-projector}
\end{equation}
Отсюда без дополнительного множителя удвоения получается
\begin{equation}
 \Tr\!\left[P_{\rm phys}(I_4\otimes D_F^4)\right]
 =2\Tr D_F^4=92,
 \qquad N_{\rm ferm}=-92.
 \label{eq:v8-full-carrier-fermion-multiplicity}
\end{equation}

BV-проверка даёт отрицательный результат для прежнего скалярного блока.
Пусть $X_{\rm orb}$ переводит двенадцать калибровочных генераторов в
координаты 30-мерного вещественного transfer-кадра. Тогда точно
\begin{equation}
 \rank X_{\rm orb}=3,
 \qquad
 \rank\!\left(X_{\rm orb}^{\mathsf T}H_{\Phi}X_{\rm orb}\right)=3,
 \qquad
 \Tr\!\left(X_{\rm orb}^{\mathsf T}H_{\Phi}X_{\rm orb}\right)=34.
 \label{eq:v8-full-carrier-goldstone-failure}
\end{equation}
Следовательно, три нарушенных калибровочных направления не лежат в ядре
$H_{\Phi}$. Это не допустимая масса голдстоунов: использованный ранее
грамов гессиан дифференцировал функционал с фиксированным фоном
инцидентности, а не калибровочно-инвариантный вакуумный функционал на
совместно движущемся фоне.

Формальная подстановка уже выведенного фермионного вклада дала бы
\begin{equation}
 N_{\rm fixed}=\frac{4659176}{3249}-92
 =\frac{4360268}{3249},
 \label{eq:v8-full-carrier-fixed-background-candidate}
\end{equation}
но это число не является физическим коэффициентом $B$, поскольку его
босонная часть не прошла голдстоуновский критерий.

\begin{theorem}[Фермионная determinant-кратность]
Уже принятую полную внешне-внутреннюю хиральную проекцию достаточно
применить к $D_F^4$. Она имеет ранг $42$ и однозначно переводит сырой
конечный след $46$ в физический фермионный четвёртый момент $92$.
\end{theorem}

\begin{nogocriterion}[Фиксированный фон не задаёт BV-гессиан]
Физический калибровочно-инвариантный вакуумный гессиан обязан зануляться
на касательных к нарушенной калибровочной орбите. Здесь орбита имеет
размерность $3$, а ограничение прежнего гессиана на неё имеет ранг $3$.
Поэтому прежний босонный ledger и число $4360268/3249$ нельзя повышать до
физического полного суперследа.
\end{nogocriterion}

\begin{protocol}
\begin{enumerate}
 \item физический хиральный носитель фермионов: \textbf{$42$};
 \item фермионный четвёртый момент: \textbf{$92$};
 \item размерность нарушенной калибровочной орбиты: \textbf{$3$};
 \item ранг прежнего гессиана на этой орбите: \textbf{$3$};
 \item голдстоуновский критерий: \textbf{не пройден};
 \item полный физический $B$: \textbf{не закрыт};
 \item следующий гейт: \textbf{реконструкция калибровочно-инвариантного
 вакуумного гессиана}.
\end{enumerate}
\end{protocol}

Машинный сертификат сохранён в
\path{s2t_v8_full_42_carrier_bv_vacuum_quotient_gate_results.json}.